\newpage

%%%%%%%%%%%%%%%%%%%%  解答（旧 prob_p67_2〜p72_1.png）  %%%%%%%%%%%%%%%%%%%%
\setlength{\columnseprule}{0.4pt}
\begin{multicols}{2}
\raggedcolumns

{\gtfamily\large 第2回　落下運動・放物運動}\par\vspace{2mm}

\Rank{A問題}
\MonN{9}{基本事項の確認}\par\nopagebreak
\begin{sisin}
軸の設定の向きに注意せよ．投げ上げる運動のときは鉛直上向きを正にとり，投げ下ろすもしくは自由落下させる場合には下向きを正にとって考えるとよい．各々の場合に，物体の加速度が正・負どちらになるのかに注意して立式せよ．
\end{sisin}
\noindent{\gtfamily 【解答】}\par\nopagebreak
\begin{komon}
\ko{1} $h=\dfrac{1}{2}\cdot 9.8\cdot 2.0^2=19.6\fallingdotseq 2.0\times 10\U{m}$\Ans\par
       $v=0+9.8\cdot 2.0=19.6\fallingdotseq 2.0\times 10\U{m/s}$\Ans
\ko{2} $v=9.8+9.8\cdot 3.0=39.2$\par
       \hspace*{1zw}$\fallingdotseq 3.9\times 10\U{m/s}$\Ans
\ko{3} $v=19.6+(-9.8)\cdot 2.0=0.0\U{m/s}$\Ans\par
       $h=0+19.6\cdot 2.0+\dfrac{1}{2}\cdot(-9.8)\cdot 2.0^2$\par
       \hspace*{1zw}$=19.6\fallingdotseq 2.0\times 10\U{m}$\Ans
\ko{4} $x=0+20\cdot 3.0=6.0\times 10\U{m}$\Ans\par
       $y=0+0\cdot 3.0+\dfrac{1}{2}\cdot 9.8\cdot 3.0^2$\par
       \hspace*{1zw}$=44.1\fallingdotseq 4.4\times 10\U{m}$\Ans
\ko{5} $v_x=20\cos 30^\circ=10\sqrt{3}$\par
       \hspace*{1zw}$\fallingdotseq 1.7\times 10\U{m/s}$\Ans\par
       $v_y=20\sin 30^\circ-9.8\cdot 1.0$\par
       \hspace*{1zw}$=2.0\times 10^{-1}\U{m/s}$\Ans
\ko{6} 最高点では速度の鉛直成分が$0\U{m/s}$なので，
       \[0=19.6\sin 30^\circ-9.8\cdot t\]
       これより，
       \[t=1.0\U{s}\Ans\]
\end{komon}

\Rank{B問題}
\MonN{10}{鉛直投げ下げと自由落下}\par\nopagebreak
\begin{sisin}
自由落下運動あるいは鉛直投げ下げ運動であるから，小石$\mathrm{A}$，$\mathrm{B}$を投げ出すところを原点$\mathrm{O}$として鉛直下向きに$y$軸をとって考えるとよい．物体の落下運動は等加速度直線運動であるから，前回学んだ三つの関係式をうまく適用して考えよ．
\end{sisin}
\noindent{\gtfamily 【解答】}\par\nopagebreak
\hspace{1zw}小石$\mathrm{A}$，$\mathrm{B}$を投げ出すところを原点$\mathrm{O}$として，鉛直下向きに$y$軸をとる．
\begin{komon}
\ko{1} $\mathrm{A}$は$\mathrm{B}$に追い付かれるまでに，$y$軸上にて初速度$0\U{m/s}$，加速度$9.8\U{m/s^2}$の等加速度直線運動を$3.0\U{s}$間だけ行っている．ゆえに，この$3.0\U{s}$間の$\mathrm{A}$の移動距離$\Delta y$は
       \[\Delta y=\frac{1}{2}\cdot 9.8\cdot 3.0^2=44.1\U{m}\Ans\]
\ko{2} $\mathrm{B}$の初速度を$v_0\U{m/s}$とおくと，$\mathrm{B}$が投げ下ろされてからの運動は，初速度$v_0$，加速度$9.8\U{m/s^2}$の等加速度直線運動だと捉えることができる．そして，この運動を$2.0\U{s}$だけ行うと$44.1\U{m}$だけ落下して$\mathrm{A}$に追い付くのだから，
       \[v_0\cdot 2.0+\frac{1}{2}\cdot 9.8\cdot 2.0^2=44.1\]
       これを解いて，
       \[v_0=12.25\fallingdotseq 1.2\times 10\U{m/s}\Ans\]
\end{komon}

\MonN{11}{鉛直投げ上げ}\par\nopagebreak
\begin{sisin}
本問は「投げ上げ」運動であるから，鉛直上向きに$y$軸をとって考えるとよい．加速度は$-g$になる．また，$y$軸の原点を地上にとる方法と屋上にとる方法との2通りが考えられるが，どちらでも構わない．ここでは地上を原点$\mathrm{O}$とした場合の解答を示す．

「鉛直投げ上げ運動の最高点では，小石の速度が$0\U{m/s}$」に着目する．
\end{sisin}
\noindent{\gtfamily 【解答】}\par\nopagebreak
\hspace{1zw}地面を原点として鉛直上向きに$y$軸をとり，小石を投げ上げた時刻を$0\U{s}$とする．
\begin{komon}
\ko{1} 小石は$y$軸上を初期位置$h$，初速度$v_0$，加速度$-g$の等加速度直線運動するので，時刻$t\U{s}$における小石の速度$v(t)\U{m/s}$および座標$y(t)\U{m}$は，
       \[v(t)=v_0+(-g)t\Eq{1}\]
       \[y(t)=h+v_0t+\frac{1}{2}(-g)t^2\Eq{2}\]
       \hspace{1zw}最高点で小石の速度は$0\U{m/s}$となるから，最高点に達する時刻$t_0\U{s}$は\MARU{1}より，
       \[0=v_0-gt_0\]
       これより，$t_0=\dfrac{v_0}{g}\U{s}$\Ans\Eq{3}\par
       \hspace{1zw}また地上から最高点までの高さは，\MARU{2}式に$t=t_0$を代入して，
       \[h+\frac{v_0^2}{2g}\U{m}\Ans\]
\ko{2} 地面に達するとき小石の座標は$y=0\U{m}$となるので，\MARU{2}より，
       \[0=h+v_0t+\frac{1}{2}(-g)t^2\]
       解の公式より，$t=\dfrac{v_0\pm\sqrt{v_0^2+2gh}}{g}$\par
       求める時刻は正であり，
       \[t=\frac{v_0+\sqrt{v_0^2+2gh}}{g}\U{s}\Ans\Eq{4}\]
       またそのときの速さは，\MARU{1}\MARU{4}より，
       \[|v|=\sqrt{v_0^2+2gh}\U{m/s}\Ans\]
\end{komon}
\begin{itembox}[l]{【参考】}
\hspace{1zw}(1)の高さでは，$2a\Delta x=v^2-v_0^2$を利用してもよい．最高点の高さを$h'$とすると，上昇距離は$h'-h$であるから，
\[2\cdot(-g)\cdot(h'-h)=0^2-v_0^2\]
からも答えが求められる．本問の座標において加速度は$-g$なのだから，
\[2\cdot g\cdot(h'-h)=v_0^2-0^2\]
\[2\cdot g\cdot(h-h')=0^2-v_0^2\]
などと立式しないように注意せよ．

\hspace{1zw}なお，(2)の2解のうち$t<0$の解は，この時刻に地上から鉛直上向きに小石を投げ上げて，時刻$t=0$に屋上を上向き速さ$v_0\U{m/s}$で通過する場合に対応する解である．
\end{itembox}

\MonN{12}{水平投射}\par\nopagebreak
\begin{sisin}
2次元面内で物体が運動する場合には，物体の運動を水平方向と鉛直方向に分解して考えるのが基本．特に本問の場合は加速度を持つのは鉛直方向のみであるから，分解して考えれば水平方向については等速度運動，鉛直方向には等加速度運動（自由落下）となり，それぞれについて直線運動と同様に扱えばよくなる．

座標軸の設定としては，本問の場合は小石を投げ出したところを原点$\mathrm{O}$とし，水平右向きに$x$軸，鉛直下向きに$y$軸をとると計算が楽になる．
\end{sisin}
\noindent{\gtfamily 【解答】}\par\nopagebreak
\hspace{1zw}小石を投げ出したところを原点$\mathrm{O}$とし，水平右向きに$x$軸，鉛直下向きに$y$軸をとる．また，投げ出した時刻を$0\U{s}$とする．小石は$x$軸方向には速度$v_0\U{m/s}$の等速度運動を，$y$軸方向には初速度$0\U{m/s}$，加速度$g\U{m/s^2}$の等加速度運動（自由落下）をするので，時刻$t\U{s}$における小石の速度の$x,y$成分，および$x,y$座標は，
\[v_x(t)=v_0\quad v_y(t)=gt\Eq{1}\]
\[x(t)=v_0t\quad y(t)=\frac{1}{2}gt^2\Eq{2}\]
と表される．
\begin{komon}
\ko{1} 地面に達するときの$y$座標は$y=h\U{m}$であるから，地面に達する時刻を$t=t_0\U{s}$とすると，\MARU{2}より，
       \[\frac{1}{2}g{t_0}^2=h\]
       よって，$t_0=\sqrt{\dfrac{2h}{g}}\U{s}$\Ans\Eq{3}
\ko{2} この間の水平方向の運動に着目すると，水平到達距離は\MARU{2}\MARU{3}より，
       \[x(t_0)=v_0\sqrt{\frac{2h}{g}}\U{m}\Ans\Eq{4}\]
\ko{3} 地面に達する直前の速度の水平成分および鉛直成分は，\MARU{1}\MARU{3}より，\par
       \hspace{1zw}水平成分：$v_x=v_0$\par
       \hspace{1zw}鉛直成分：$v_y=gt_0=\sqrt{2gh}$\par
       \hspace{1zw}よってこのときの速さ$|v|$（速度ベクトルの大きさ）は，
       \[|v|=\sqrt{{v_x}^2+{v_y}^2}=\sqrt{v_0^2+2gh}\U{m/s}\Ans\]
\end{komon}

\MonN{13}{斜方投射}\par\nopagebreak
\begin{sisin}
本問も2次元面内での物体の運動なので，$x,y$各軸方向の運動に分けて考えればよいが，斜方投射の場合にはボールの鉛直方向の運動が鉛直投げ上げ運動になっていることに着目すると素早く解くことができる．

［11］同様，鉛直上向きを$y$軸の正とするので，重力加速度は$-g$と考える．
\end{sisin}
\noindent{\gtfamily 【解答】}\par\nopagebreak
\hspace{1zw}ボールを投げ出した時刻を$0\U{s}$とする．
\begin{komon}
\ko{1} ボールは$x$軸方向には速度$v_0\cos\theta\U{m/s}$の等速度運動を，$y$軸方向には初速度$v_0\sin\theta\U{m/s}$，加速度$-g\U{m/s^2}$の鉛直投げ上げ運動をする．よって，時刻$t\U{s}$におけるボールの速度の$x,y$成分は，
       \[v_x(t)=v_0\cos\theta\]
       \[v_y(t)=v_0\sin\theta+(-g)t\Eq{1}\]
       \hspace{1zw}また，物体の$x,y$座標は，
       \[x(t)=v_0\cos\theta\cdot t\U{m}\]
       \[y(t)=v_0\sin\theta\cdot t+\frac{1}{2}(-g)t^2\U{m}\Ans\Eq{2}\]
\ko{2} 最高点ではボールの速度の鉛直成分が$0\U{m/s}$であるから，最高点に達する時刻を$t_1\U{s}$とおくと\MARU{1}より，
       \[0=v_0\sin\theta-gt_1\]
       これより，$t_1=\dfrac{v_0\sin\theta}{g}\U{s}$\Ans\Eq{3}\par
       \hspace{1zw}またこのときの$y$座標は，\MARU{2}\MARU{3}より，
       \[y(t_1)=\frac{v_0^2\sin^2\theta}{2g}\U{m}\Ans\]
\end{komon}
\begin{bekkai}
\hspace{1zw}$y$軸方向の等加速度運動に着目すると，最高点での速度の鉛直成分は$0\U{m/s}$なので，高さを$h$として，
\[2\cdot(-g)\cdot h=0^2-(v_0\sin\theta)^2\]
が成立する．これより，
\[h=\frac{v_0^2\sin^2\theta}{2g}\U{m}\Ans\]
\end{bekkai}
\begin{komon}
\ko{3} 再び地面に戻ってきたときの時刻を$t_2\U{s}$とおくと，$y$座標が$0\U{m}$となるので，\MARU{2}より，
       \[0=v_0\sin\theta\cdot t_2+\frac{1}{2}\cdot(-g){t_2}^2\]
       が成り立つ．ボールを投げ出した時刻が$0\U{s}$であり，$t_2>0$であるから，
       \[t_2=\frac{2v_0\sin\theta}{g}\U{s}\]
       \hspace{1zw}よって水平到達距離は$x$軸方向に着目して，\MARU{2}より，
       \[x(t_2)=\frac{2v_0^2\sin\theta\cos\theta}{g}\U{m}\Ans\]
\end{komon}

\MonN{14}{斜方投射}\par\nopagebreak
\begin{sisin}
平面内の運動なので，水平方向と鉛直方向に分解する．初速が斜め右上向きであるから，水平右向きに$x$軸を，鉛直上向きに$y$軸をとって考えよう．加速度は$-g\U{m/s^2}$になることに注意．
\end{sisin}
\noindent{\gtfamily 【解答】}\par\nopagebreak
\hspace{1zw}ボールを投げた時刻を$0\U{s}$とする．ボールは$x$軸方向には速度$v_0\cos\theta\U{m/s}$の等速度運動を，$y$軸方向には初速度$v_0\sin\theta\U{m/s}$，加速度$-g\U{m/s^2}$の等加速度運動をする．よって，時刻$t\U{s}$におけるボールの速度の$x,y$成分$v_x(t),v_y(t)\U{m/s}$は，
\[v_x(t)=v_0\cos\theta\]
\[v_y(t)=v_0\sin\theta+(-g)t\Eq{1}\]
と表され，ボールの$x,y$座標$x(t),y(t)\U{m}$は，
\[x(t)=v_0\cos\theta\cdot t\U{m}\Eq{2}\]
\[y(t)=v_0\sin\theta\cdot t+\frac{1}{2}(-g)t^2\U{m}\]
と表される．
\begin{komon}
\ko{1} 最高点では，ボールの速度の鉛直成分が$0\U{m/s}$であるから，\MARU{1}より，
       \[0=v_0\sin\theta-gt\]
       これより，$t=\dfrac{v_0\sin\theta}{g}\U{s}$\Ans
\ko{2} 地面に達するとき，すなわちボールの$y$座標が$-h\U{m}$となる時刻を$t=t_0\U{s}$とおくと，
       \[-h=v_0\sin\theta\cdot t_0-\frac{1}{2}g{t_0}^2\]
       であるから，
       \[t_0=\frac{v_0\sin\theta\pm\sqrt{v_0^2\sin^2\theta+2gh}}{g}\]
       $t_0>0$であるから，地面に達するまでの時間は，
       \[\frac{v_0\sin\theta+\sqrt{v_0^2\sin^2\theta+2gh}}{g}\U{s}\]
       \hspace*{\fill}\Ans\hspace{1zw}\mbox{$\cdots$\MARU{3}}
\ko{3} このときの水平到達距離は，\MARU{2}\MARU{3}より，
       \[x(t_0)=\frac{v_0^2\cos\theta\sin\theta}{g}\]
       \[\hspace*{2zw}+\frac{v_0\cos\theta\sqrt{v_0^2\sin^2\theta+2gh}}{g}\U{m}\Ans\]
\ko{4} 地面に達する直前のボールの速度の$x,y$成分は，\MARU{1}\MARU{3}より，\par
       \hspace{1zw}$x$成分：$v_x(t_0)=v_0\cos\theta\U{m/s}$\par
       \hspace{1zw}$y$成分：$v_y(t_0)=-\sqrt{v_0^2\sin^2\theta+2gh}\U{m/s}$\par
       \hspace*{\fill}\Ans
\end{komon}

\MonN{15}{自由落下と鉛直投げ上げ}\par\nopagebreak
\begin{sisin}
運動は一直線上で起こるから，鉛直方向に$y$軸をとって考えればよいが，上向き・下向きのどちらを正とするかで2通りある．もちろんどちらを用いてもよいが，本問の場合は地上を原点にとるのが楽そうであるから，地上を原点にとって鉛直上向きに$y$軸をとればよいだろう．鉛直上向きを正としたのなら，速度・加速度は鉛直上向きの場合を正，鉛直下向きの場合を負として式に代入しなければならない．加速度は$-g\U{m/s^2}$を用いなければならないことにも注意せよ．
\end{sisin}
\noindent{\gtfamily 【解答】}\par\nopagebreak
\hspace{1zw}地面を原点$\mathrm{O}$として鉛直上向きに$y$軸をとり，$\mathrm{A}$と$\mathrm{B}$を投げ出した時刻を$0\U{s}$とする．
\begin{komon}
\ko{1} $\mathrm{A}$は$y=h$の初期位置から初速度$0\U{m/s}$，加速度$-g\U{m/s^2}$の等加速度直線運動を，また$\mathrm{B}$は$y=0\U{m}$の初期位置から初速度$+v_0\U{m/s}$，加速度$-g\U{m/s^2}$の等加速度直線運動をするので，時刻$t\U{s}$における$\mathrm{A}$，$\mathrm{B}$の速度および$y$座標は，
       \[\mathrm{A}:v_{\mathrm{A}}(t)=0+(-g)\cdot t\]
       \[y_{\mathrm{A}}(t)=h+0\cdot t+\frac{1}{2}(-g)\cdot t^2\Eq{1}\]
       \[\mathrm{B}:v_{\mathrm{B}}(t)=v_0+(-g)\cdot t\]
       \[y_{\mathrm{B}}(t)=0+v_0t+\frac{1}{2}(-g)\cdot t^2\Eq{2}\]
       \hspace{1zw}$\mathrm{A}$と$\mathrm{B}$が衝突するときとは，$y$座標が一致するときなので，この時刻を$t=t_0\U{s}$とすると，\MARU{1}\MARU{2}より，
       \[h-\frac{1}{2}g{t_0}^2=v_0t_0-\frac{1}{2}g{t_0}^2\]
       % 誤植訂正（2026-08-29）：元の問題集はこの式が h = \frac{1}{2}gt_0^2 = v_0t_0 - \frac{1}{2}gt_0^2 と
       %   なっていた（最初の等号が誤り）．正しくは左辺が y_A(t_0) すなわち h-\frac{1}{2}gt_0^2．
       %   次行の t_0 = h/v_0 とも整合する．
       これより，$t_0=\dfrac{h}{v_0}\U{s}$\Ans\Eq{3}
\ko{2} このときの地面からの高さは，\MARU{1}$\sim$\MARU{3}より，
       \[y_{\mathrm{A}}(t_0)=y_{\mathrm{B}}(t_0)=h-\frac{gh^2}{2v_0^2}\U{m}\Ans\]
\end{komon}

\Rank{C問題}
\MonN{16}{鉛直投げ上げと等速度運動}\par\nopagebreak
\begin{sisin}
本問を捉えるには「いつの時刻に，気球および小石がどこにいるのか」をまずはっきりさせる必要がある．そのためには，小石を投げた瞬間，その$3.0\U{s}$後，さらに$4.0\U{s}$後の三つの瞬間について，小石，気球，地面の三者の位置関係を図示し，すぐに分かる部分について距離を記入せよ．例えば，最初の瞬間と$3.0\U{s}$後，さらに$4.0\U{s}$後の気球の位置関係は，最初の位置に対してそれぞれ下方向に$21\U{m}$，$49\U{m}$だけずれている．このような図を図示することがまず第一歩である．

各物体の座標を座標軸上に置き，ひたすら数式計算で考えていく，というのも一つの方法としてはあるが，その場合でも上に述べたような簡単な図を描くだけでケアレスミスが防げる．
\end{sisin}
\noindent{\gtfamily 【解答】}\par\nopagebreak
\hspace{1zw}小石を投げ出した地点を原点$\mathrm{O}$として鉛直上向きに$y$軸をとり，小石を投げ出した時刻を$0\U{s}$とする．小石の地面に対する鉛直上向きの初速度を$v_0\U{m/s}$とすると，時刻$t\U{s}$における小石の速度$v(t)\U{m/s}$および座標$y(t)\U{m}$は，
\[\left\{\begin{array}{l}
  v(t)=v_0+(-9.8)t\\[1mm]
  y(t)=v_0t+\dfrac{1}{2}(-9.8)t^2
\end{array}\right.\Eq{1}\]
である．また，気球は一定速度$-7.0\U{m/s}$で運動しつづけるので，時刻$t\U{s}$における気球の速度および座標は，
\[V(t)=-7.0\U{m/s},\ \ Y(t)=-7.0\cdot t\U{m}\Eq{2}\]
となる．
\begin{komon}
\ko{1} 小石が気球を追い越す時刻$3.0\U{s}$における気球の位置は$Y(3.0)=-21$であるから，気球の落下距離は，
       \[|Y|=2.1\times 10\U{m}\Ans\Eq{3}\]
\ko{2} 時刻$3.0\U{s}$では小石と気球の$y$座標が等しくなるので，\MARU{1}\MARU{3}より
       \[v_0\cdot 3.0-\frac{1}{2}\cdot 9.8\cdot 3.0^2=-21.0\]
       よって，$v_0=7.7\U{m/s}$\Eq{4}\par
       \hspace{1zw}これより，時刻$3.0\U{s}$での小石の速さは，
       \[|v(3.0)|=21.7\fallingdotseq 2.2\times 10\U{m/s}\Ans\]
\ko{3} 小石は時刻$t=3.0+4.0=7.0\U{s}$に地上に達するので，地上は$y(7.0)\U{m}$の位置にある．この時刻における気球の座標は$Y(7.0)\U{m}$であり，これは地上からは$Y(7.0)-y(7.0)\U{m}$の高さである．
       \[Y(t)-y(t)=-(7.0+v_0)t-\frac{1}{2}(-9.8)t^2\]
       % 誤植訂正（2026-08-29）：元の問題集は右辺第2項が +\frac{1}{2}(-9.8)t^2 だったが，
       %   Y(t)-y(t) = -7.0t-\{v_0t+\frac{1}{2}(-9.8)t^2\} なので符号は -\frac{1}{2}(-9.8)t^2 が正しい．
       %   これで Y(7.0)-y(7.0)=137.2 となり次行の答と整合する（誤植のままだと -343 になる）．
       であるから，求める高さは，
       \[Y(7.0)-y(7.0)=137.2\fallingdotseq 1.4\times 10^2\U{m}\Ans\]
\end{komon}
\begin{itembox}[l]{【参考】}
\hspace{1zw}上述の解答では，座標軸の原点を最初の点としたが，地上を原点とするのも一つの手である．この場合は，最初の気球の地上からの高さを$h$として考えればよい．
\end{itembox}

\MonN{17}{放物運動}\par\nopagebreak
\begin{sisin}
台車からボールを真上に打ち上げた，とは，台車に対して真上に打ち上げた，という意味である．すなわち，台車に対するボールの相対速度が鉛直上向きであったことを意味する．よって，地面から見た場合には，ボールは鉛直上向きに$u\U{m/s}$を持つとともに，水平右向きに台車の速度$v\U{m/s}$を持つ．ゆえにこのボールは地面から見て放物線を描いて$\mathrm{B}$点に達することになる．
\end{sisin}
\noindent{\gtfamily 【解答】}\par\nopagebreak
\begin{komon}
\ko{1} ボールは投げ出された直後，地面から見て鉛直上向きに$u\U{m/s}$，水平右向きに$v\U{m/s}$の速度を持つ．鉛直方向の投げ上げ運動に着目すると上昇に$\dfrac{u}{g}\U{s}$，下降に$\dfrac{u}{g}\U{s}$を要するので，$\mathrm{A}\to\mathrm{B}$に要する時間は$\dfrac{2u}{g}\U{s}$である．\par
       \hspace{1zw}よってこの間にボールが水平方向に進んだ距離，すなわち$\mathrm{AB}$間の距離は，
       \[\overline{\mathrm{AB}}=v\cdot\frac{2u}{g}\U{m}\Ans\]
       となる．\par
       \hspace{1zw}またボールが達した最高点の高さ$h$は，鉛直方向の運動に着目して，
       \[2\cdot(-g)\cdot h=0^2-u^2\]
       よって，$h=\dfrac{u^2}{2g}\U{m}$\Ans
\end{komon}
\begin{itembox}[l]{【参考】}
\hspace{1zw}この間に台車が進んだ距離について考えると，これも全く同じく$v\cdot\dfrac{2u}{g}\U{m}$となるので，ボールは地上に戻ってきたときにちょうど台車に衝突することになる．
\end{itembox}
\begin{komon}
\ko{2} 台車に対するボールの鉛直上向き投げ上げ速度を$u'$とすると，この場合のボールの水平到達距離は，(1)の場合の答えを$v\to\dfrac{1}{2}v$，$u\to u'$と置き換えることにより，
       \[\frac{1}{2}v\cdot\frac{2u'}{g}\U{m}\]
       となる．これがちょうど$\mathrm{AB}$間の距離$\overline{\mathrm{AB}}=v\cdot\dfrac{2u}{g}\U{m}$に等しくなるので，
       \[\frac{1}{2}v\cdot\frac{2u'}{g}=v\cdot\frac{2u}{g}\U{m}\]
       よって，$\dfrac{u'}{u}=2$\Ans
\ko{3} このときのボールの最高点までの高さを$h'$とすると，これは(1)の場合の$h$の式において$u\to u'$と置き換えればすぐに求められ，
       \[h'=\frac{u'^2}{2g}\U{m}\]
       ゆえに，
       \[\frac{h'}{h}=4\Ans\]
       % 誤植訂正（2026-08-29）：元の問題集はここが \frac{h'}{u}=4 だったが，
       %   設問は「最高点の高さは先ほどの場合の何倍か」なので h'/h が正しい．
       %   h'/h=(u'/u)^2=2^2=4 で値も整合する．
\end{komon}

\MonN{18}{モンキーハンティング}\par\nopagebreak
\begin{sisin}
本問は，落下する物体をライフルで打ち落とす場合に，どの方向を狙って撃てばよいのかを示す問題である．物体$\mathrm{P}$が落下してくることを考慮すると，$\mathrm{A}$点よりも下方を狙った方がよさそうに思えるが，実際には$\mathrm{A}$点を狙って発射すれば落下してきた$\mathrm{P}$に当たる．
\end{sisin}
\noindent{\gtfamily 【解答】}\par\nopagebreak
\begin{komon}
\ko{1} $\mathrm{P}$は自由落下運動をするので，時刻$t\U{s}$における$\mathrm{P}$の速度の$x,y$成分および$x,y$座標は，
       \[V_x(t)=0\quad V_y(t)=-gt\Eq{1}\]
       \[X(t)=l\quad Y(t)=h+\frac{1}{2}(-g)t^2\Eq{2}\]
       \hspace{1zw}$\mathrm{Q}$は放物運動をするので，時刻$t\U{s}$における$\mathrm{Q}$の速度の$x,y$成分および$x,y$座標は，
       \[\left\{\begin{array}{l}
         v_x(t)=v_0\cos\theta\\[1mm]
         v_y(t)=v_0\sin\theta+(-g)t
       \end{array}\right.\Eq{3}\]
       \[\left\{\begin{array}{l}
         x(t)=v_0\cos\theta\cdot t\\[1mm]
         y(t)=v_0\sin\theta\cdot t+\dfrac{1}{2}(-g)t^2
       \end{array}\right.\Eq{4}\]
       \hspace{1zw}よって，時刻$t\U{s}$における$\overline{\mathrm{PQ}}$は，\MARU{2}\MARU{4}より，
       \[\overline{\mathrm{PQ}}=\sqrt{(v_0\cos\theta\cdot t-l)^2+(v_0\sin\theta\cdot t-h)^2}\]
       \hspace*{\fill}$\U{m}$\Ans
\ko{2} 時刻$t\U{s}$における$\mathrm{P}$から見た$\mathrm{Q}$の相対速度は，
       \[\overrightarrow{v_{\mathrm{PQ}}}=\overrightarrow{v_{\mathrm{Q}}}-\overrightarrow{v_{\mathrm{P}}}\Eq{5}\]
       よって，\MARU{1}\MARU{3}\MARU{5}より，
       \[\overrightarrow{v_{\mathrm{PQ}}}=\left(\begin{array}{c}
         v_0\cos\theta\\ v_0\sin\theta
       \end{array}\right)\]
       であり，
       \[\left\{\begin{array}{l}
         \mbox{水平成分：}v_0\cos\theta\U{m/s}\\[1mm]
         \mbox{鉛直成分：}v_0\sin\theta\U{m/s}
       \end{array}\right.\Ans\]
\ko{3} $\mathrm{P}$と$\mathrm{Q}$が衝突するのは$x,y$座標が同時刻に等しくなるときである．衝突時刻を$t=t_0$とすると，$x$座標が等しくなるのは，\MARU{2}\MARU{4}より，
       \[v_0\cos\theta\cdot t_0=l\]
       よって，$t_0=\dfrac{l}{v_0\cos\theta}$\Eq{6}\par
       \hspace{1zw}この時刻に$y$座標も等しくなればよい．\MARU{2}\MARU{4}より，
       \[v_0\sin\theta\cdot t_0-\frac{1}{2}g{t_0}^2=h-\frac{1}{2}g{t_0}^2\]
       これを解いて，
       \[h=l\tan\theta\U{m}\Ans\Eq{7}\]
\ko{4} $\mathrm{P}$と$\mathrm{Q}$が衝突する$y$座標は，\MARU{2}\MARU{6}\MARU{7}より，
       \[y=h-\frac{1}{2}g\left(\frac{l}{v_0\cos\theta}\right)^2\Eq{8}\]
       これが$y>0$であればよいので，
       \[v_0>\frac{l}{\cos\theta}\sqrt{\frac{g}{2h}}=\sqrt{\frac{gl}{2\sin\theta\cos\theta}}\U{m/s}\]
       \hspace*{\fill}\Ans
\end{komon}
\begin{itembox}[l]{【参考】}
\hspace{1zw}(2)の結果から分かるように，物体$\mathrm{P}$の上から見ると物体$\mathrm{Q}$は初速度を保ったまま$\mathrm{P}$に接近してくるように見える．ゆえに$\mathrm{Q}$が$\mathrm{P}$に衝突する条件は，$\mathrm{P}$から見たときに$\mathrm{Q}$が直進しつづければ$\mathrm{P}$に当たるようになっていることであり，最初の瞬間において$\mathrm{A}$点を狙って撃てばよいことが分かる．
\end{itembox}

\end{multicols}
